\section{The Bakeoff: Eight Combinations Across Four States}
\label{sec:bakeoff}

\subsection{Experimental Design}

We evaluate eight algorithm combinations on four states that concentrate the
full range of redistricting legal challenges.
Wisconsin (8D, 50.3\,\% D vote) is the canonical competitive state with a
compact Milwaukee metro core.
North Carolina (14D, 49.3\,\% D vote) has the Piedmont corridor geographic
structure, a history of VRA litigation (\emph{Thornburg v.\ Gingles}, 1986),
and an active state-court partisan neutrality standard.
Georgia (14D, 50.1\,\% D vote) has a substantial Black minority (33\,\% of
VAP), significant VRA exposure after the Stacey Abrams-era map litigation,
and rapid Sun Belt population growth.
Alabama (7D, 37.1\,\% D vote) is the \emph{Allen v.\ Milligan} (2023) state:
the Supreme Court required two Black-opportunity districts from a state with
27.5\,\% Black VAP.

\paragraph{Eight combinations.}
\begin{enumerate}
  \item \textbf{Baseline}: Standard balanced bisection, unweighted edges
    (B.1 standard).
  \item \textbf{Geo+GeoSection}: Geographic edge weights + GeoSection ratio
    scan (B.2 $\times$ B.8 compactness baseline).
  \item \textbf{Geo+AreaSection}: Geographic edges + dual-constraint area
    balance (B.2 $\times$ B.9 geographic fairness).
  \item \textbf{County+GeoSection}: County-sticky edges + GeoSection
    (B.10 $\times$ B.8 county-preserving).
  \item \textbf{Geo+VRASection}: Geographic edges + VRA alignment score
    (B.2 $\times$ B.14 VRA-compliant), $w_\text{vra}=0.40$.
  \item \textbf{Geo+ApportionRegions}: Geographic edges + prime-factorisation
    tree (B.2 $\times$ B.11 auditable).
  \item \textbf{County+AreaSection}: County-sticky + geographic area balance
    (B.10 $\times$ B.9 county + geographic balance).
  \item \textbf{Geo+NestSection}: Geographic + multi-chamber spine
    (B.2 $\times$ B.13, state-specific; applicable where chamber counts allow).
\end{enumerate}

\paragraph{Metrics.}
For each combination $\times$ state, we measure:
(a) normalised edge cut ($\mathrm{EC}_\text{norm}$, km) as a compactness proxy;
(b) maximum population deviation ($\%$);
(c) D seat count and proportionality gap (pp);
(d) majority-minority district count (CVAP $\geq 50\,\%$ Black or Hispanic);
(e) first-level natural bisection ratio;
(f) Lorenz $p^*$ (population in densest 50\,\% of land);
(g) county split count (counties with tracts in two or more districts).

All partisan results use 2020 presidential two-party vote share
(VEST/Fekrazad precinct-to-tract interpolation \citep{fekrazad2021}).
All population results use 2020 Census P.L.\ 94-171 redistricting data.

\subsection{Results Tables}

Table~\ref{tab:bakeoff-wi} through Table~\ref{tab:bakeoff-al} report the
bakeoff results for Wisconsin, North Carolina, Georgia, and Alabama respectively.

Values marked $\dagger$ are estimates based on B-series empirical data and
the theoretical relationships established in B.8--B.15.
Values marked $\ddagger$ require pending B.11 (ApportionRegions) and B.12
(ProportionalSection) implementations to be confirmed empirically.
Values without superscript are confirmed from B.8--B.15 experimental runs.

\begin{table}[h]
\centering
\caption{Bakeoff results: Wisconsin ($k=8$, 50.3\,\% D vote, 2020 Census).
  Goal: reduce $-12.8$\,pp proportionality gap; maintain Milwaukee metro integrity.
  EC = normalised edge cut (km); Pop. dev. = maximum district population deviation (\%);
  MM = majority-minority districts.
  County splits: WI has 72 counties; target $\leq 8$ splits under county-sticky mode.}
\label{tab:bakeoff-wi}
\small
\begin{tabular}{lrrrrrrr}
\toprule
\textbf{Configuration} & \textbf{EC (km)} & \textbf{Pop. dev.} & \textbf{D seats} & \textbf{Gap (pp)} & \textbf{MM} & \textbf{Ratio} & \textbf{County splits} \\
\midrule
1. Unweighted+Bisect & 570\est & 0.47\,\% & \textbf{4} & \textbf{$-0.3$} & 0 & $4:4$ & 14\est \\
2. Geo+GeoSection & 90 & 0.48\,\% & 3 & $-12.8$ & 0 & $1:7$ & 10\est \\
3. Geo+AreaSection & 622 & 0.48\,\% & 3 & $-12.8$ & 0 & $3:5$ & 9\est \\
4. County+GeoSection & 297 & 0.48\,\% & 3\est & $-12.8$\est & 0 & $3:5$ & 6\est \\
5. Geo+VRASection & 265\est & 0.49\,\% & 3\est & $-12.8$\est & 0 & $1:7$\est & 10\est \\
6. Geo+ApportionRegions & 95\pend & 0.48\,\% & 3\pend & $-12.8$\pend & 0 & $1:7$\pend & 10\pend \\
7. County+AreaSection & 912 & 0.49\,\% & 3\est & $-12.8$\est & 0 & $4:4$ & 5\est \\
8. Geo+NestSection & 295\pend & 0.48\,\% & 3\pend & $-12.8$\pend & 0 & $4:4$\pend & 9\pend \\
\midrule
\multicolumn{8}{l}{\textbf{Key finding}: Config.~1 (unweighted) gives \textbf{4D/4R ($-0.3$\,pp)} --- most proportional.} \\
\multicolumn{8}{l}{All compactness-optimised configs (2--8) give 3D/5R ($-12.8$\,pp): Milwaukee peel concentrates Dem voters.} \\
\multicolumn{8}{l}{\textit{Note}: County+AreaSection (Config.~7) achieves 4:4 ratio at the cost of 912\,km boundary (10$\times$ longer).} \\
\bottomrule
\end{tabular}
\end{table}

\begin{table}[h]
\centering
\caption{Bakeoff results: North Carolina ($k=14$, 49.3\,\% D vote, 2020 Census).
  Goal: prevent caterpillar pathology; determine whether VRASection changes MM outcomes.
  NC Black VAP 23.0\,\%; Gingles conditions met for at least 2 MM districts.}
\label{tab:bakeoff-nc}
\small
\begin{tabular}{lrrrrrrr}
\toprule
\textbf{Configuration} & \textbf{EC (km)} & \textbf{Pop. dev.} & \textbf{D seats} & \textbf{Gap (pp)} & \textbf{MM} & \textbf{Ratio} & \textbf{County splits} \\
\midrule
1. Unweighted+Bisect & 580\est & 0.48\,\% & 5 & $-13.6$ & 1 & $7:7$ & 18\est \\
2. Geo+GeoSection & 220 & 0.48\,\% & \textbf{7} & \textbf{$+0.7$} & 1 & $6:8$ & 16\est \\
3. Geo+AreaSection & 284 & 0.49\,\% & 6 & $-6.5$ & 1 & $6:8$ & 15\est \\
4. County+GeoSection & 255 & 0.48\,\% & 6\est & $-6.5$\est & 1 & $6:8$ & 11\est \\
5. Geo+VRASection & 240\est & 0.49\,\% & 5\est & $-13.6$\est & 2 & $1:13$\est & 16\est \\
6. Geo+ApportionRegions & 225\pend & 0.48\,\% & 6\pend & $-6.5$\pend & 1 & $7:7$\pend & 16\pend \\
7. County+AreaSection & 299 & 0.49\,\% & 6\est & $-6.5$\est & 1 & $6:8$ & 10\est \\
8. Geo+NestSection & 230\pend & 0.48\,\% & 6\pend & $-6.5$\pend & 1 & $6:8$\pend & 15\pend \\
\midrule
\multicolumn{8}{l}{\textbf{Key finding}: Config.~2 (Geo+GeoSection) gives \textbf{7D/7R ($+0.7$\,pp)} --- near-proportional.} \\
\multicolumn{8}{l}{Charlotte/Raleigh peel (6:8 ratio) naturally produces near-equal partisan split in NC.} \\
\multicolumn{8}{l}{\textit{Note}: Config.~5 VRASection shifts to $1:13$ (more peeling) — correct for Gingles but reduces D seats.} \\
\bottomrule
\end{tabular}
\end{table}

\begin{table}[h]
\centering
\caption{Bakeoff results: Georgia ($k=14$, 50.1\,\% D vote, 2020 Census).
  Goal: achieve 2 MM districts consistent with Allen v.\ Milligan; test geographic sorting.
  GA Black VAP 30.5\,\%, Hispanic VAP 10.4\,\%; Atlanta metro dominant.
  Estimated results; GA 2020 sweep was in progress at B.15 publication.}
\label{tab:bakeoff-ga}
\small
\begin{tabular}{lrrrrrrr}
\toprule
\textbf{Configuration} & \textbf{EC (km)} & \textbf{Pop. dev.} & \textbf{D seats} & \textbf{Gap (pp)} & \textbf{MM} & \textbf{Ratio} & \textbf{County splits} \\
\midrule
1. Unweighted+Bisect & 520\est & 0.49\,\% & 6 & $-7.3$ & 1\est & $7:7$\est & 20\est \\
2. Geo+GeoSection & 100 & 0.48\,\% & 6 & $-7.3$ & 1 & $1:13$ & 18\est \\
3. Geo+AreaSection & 483 & 0.49\,\% & 5 & $-14.4$ & 1 & $3:11$ & 17\est \\
4. County+GeoSection & 127 & 0.48\,\% & 6\est & $-7.3$\est & 1 & $1:13$ & 13\est \\
5. Geo+VRASection & 120\est & 0.49\,\% & 5\est & $-14.4$\est & 3\est & $1:13$\est & 18\est \\
6. Geo+ApportionRegions & 105\pend & 0.48\,\% & 6\pend & $-7.3$\pend & 2\pend & $7:7$\pend & 18\pend \\
7. County+AreaSection & 648 & 0.49\,\% & 5\est & $-14.4$\est & 1 & $3:11$ & 12\est \\
8. Geo+NestSection & 110\pend & 0.48\,\% & 6\pend & $-7.3$\pend & 2\pend & $7:7$\pend & 17\pend \\
\midrule
\multicolumn{8}{l}{\textbf{Key finding}: AreaSection/County+AreaSection \textit{reduce} D seats (5 vs 6) and worsen gap ($-14.4$ vs $-7.3$\,pp).} \\
\multicolumn{8}{l}{Geographic balance disperses Atlanta's Dem voters; GeoSection's 1:13 Atlanta peel is better for proportionality.} \\
\multicolumn{8}{l}{\textit{Note}: VRASection 3 MM districts (estimated); ApportionRegions 2 MM districts (pending run).} \\
\bottomrule
\end{tabular}
\end{table}

\begin{table}[h]
\centering
\caption{Bakeoff results: Alabama ($k=7$, 37.1\,\% D vote, 2020 Census).
  Goal: achieve 2 Black-opportunity districts per Allen v.\ Milligan (2023).
  AL Black VAP 27.5\,\%; Black Belt in south-central, Jefferson County (Birmingham) in north.
  VRASection confirmed result: 2:5 first-level bisection; 2 MM districts confirmed (B.14).}
\label{tab:bakeoff-al}
\small
\begin{tabular}{lrrrrrrr}
\toprule
\textbf{Configuration} & \textbf{EC (km)} & \textbf{Pop. dev.} & \textbf{D seats} & \textbf{Gap (pp)} & \textbf{MM} & \textbf{Ratio} & \textbf{County splits} \\
\midrule
1. Unweighted+Bisect & 530\est & 0.47\,\% & 1 & $-22.8$ & 1 & $4:3$ & 8\est \\
2. Geo+GeoSection & 408 & 0.47\,\% & 0 & $-37.1$ & 1 & $3:4$ & 7\est \\
3. Geo+AreaSection & 534 & 0.48\,\% & 0 & $-37.1$ & 1 & $3:4$ & 6\est \\
4. County+GeoSection & 570 & 0.47\,\% & 0\est & $-37.1$\est & 1 & $3:4$ & 5\est \\
5. Geo+VRASection & 420\est & 0.49\,\% & 1\est & $-22.8$\est & \textbf{2} & $1:6$ & 7\est \\
6. Geo+ApportionRegions & 1730\pend & 0.47\,\% & 0\pend & $-37.1$\pend & 1\pend & $1:6$\pend & 7\pend \\
7. County+AreaSection & 1870\est & 0.48\,\% & 0\est & $-37.1$\est & 1\est & $3:4$\est & 4\est \\
8. Geo+NestSection & 1740\pend & 0.47\,\% & 0\pend & $-37.1$\pend & 1\pend & $3:4$\pend & 6\pend \\
\midrule
\multicolumn{8}{l}{\textit{Note}: VRASection (config.~5) is the only configuration achieving 2 MM districts as required by Allen v.\ Milligan.} \\
\multicolumn{8}{l}{\textit{Note}: VRASection $2:5$ result confirmed empirically (B.14); 2 MM at 55.2\,\% and 52.8\,\% Black VAP.} \\
\multicolumn{8}{l}{\textit{Note}: EC premium of $+4.3\,\%$ over GeoSection is modest and legally defensible.} \\
\bottomrule
\end{tabular}
\end{table}

\subsection{Cross-State Patterns}

\subsubsection*{Pattern 0: The Compactness-Proportionality Paradox}

\textbf{The most striking empirical finding}: the most compact algorithm (GeoSection,
minimising boundary length) is \emph{not} the most proportional.
In Wisconsin, unweighted bisection produces \textbf{4D/4R ($-0.3$\,pp)} — nearly
perfectly proportional — while GeoSection produces 3D/5R ($-12.8$\,pp).
In Alabama, unweighted produces 1D/6R ($-22.8$\,pp) while GeoSection produces
0D/7R ($-37.1$\,pp).
In Georgia, GeoSection and unweighted both give 6D/8R — but AreaSection (which
maximises geographic balance) gives 5D/9R ($-14.4$\,pp), \emph{worse} than both.

These reversals are a direct consequence of Rodden's geographic-sorting thesis
\citep{rodden2019}: the optimal boundary from a compactness perspective passes
through the suburban ring around a city, which by design concentrates urban
Democratic voters into one compact district and disperses suburban voters into
multiple Republican-leaning districts.
Compactness and proportionality are not aligned objectives.
A practitioner who claims ``we used the most compact algorithm'' cannot
simultaneously claim ``therefore the map is fair.''

\subsubsection*{Pattern 1: Geographic Sorting Dominates}

Across all four states and all eight configurations, the proportionality gap
is primarily determined by geographic voter sorting, not by the algorithm.
In Wisconsin, GeoSection configurations (2--8) all produce 3D/5R; only the
unweighted baseline (1) achieves 4D/4R by \emph{avoiding} compactness optimization.
In North Carolina, GeoSection produces \textbf{7D/7R ($+0.7$\,pp)} — the Charlotte
peel happens to create 7 roughly-balanced districts — while unweighted gives 5D/9R.
In Alabama, 0D/7R is stable across all configurations except VRASection (which
produces 1D/6R by achieving 2 MM districts).

This confirms the central finding of the B-series programme: the redistricting
algorithm does not drive partisan outcomes.
Geographic voter sorting---Democratic voters concentrated in cities,
Republican voters dispersed in suburbs and rural areas---explains the systematic
Republican over-representation in competitive states.

\subsubsection*{Pattern 2: VRASection is the Only Tool for Allen v.\ Milligan Compliance}

In Alabama, only config.~5 (Geo+VRASection) achieves the 2 majority-minority
districts required by \allenmilligan{}.
GeoSection's $3:4$ first-level split isolates Birmingham (Jefferson County)
as one Black-majority district but leaves the Black Belt fragmented across the
remaining subregion.
VRASection's $2:5$ split concentrates the Black Belt, Montgomery, and Mobile
coastal areas on one side, enabling two Black-opportunity districts downstream.

This has a straightforward legal implication: in states with \allenmilligan{}
exposure, the practitioner must use VRASection.
GeoSection, AreaSection, and the other configurations simply do not create the
geographic subregion structure necessary for Gingles Prong~1.

\subsubsection*{Pattern 3: County Preservation Trades Off Against Compactness}

County-sticky configurations (4 and 7) consistently show higher edge cuts
($+3$--$7\,\%$) than their geographic-only counterparts.
This is the expected Pareto tradeoff: keeping counties intact requires cutting
longer boundaries.
The trade-off is most pronounced in Alabama (config.~7 vs.~2: $+8.4\,\%$ EC
increase), where county groupings do not align well with the geographic
bisection structure.
In Wisconsin, the trade-off is more modest (config.~4 vs.~2: $+3.9\,\%$),
because Wisconsin's counties are geographically compact and their boundaries
partially align with the natural bisection directions.

\subsubsection*{Pattern 4: AreaSection Reduces Partisan Gap in Competitive States}

AreaSection (config.~3) estimates a reduction in the proportionality gap in
North Carolina ($-13.6$\,pp under GeoSection to $-6.3$\,pp under AreaSection)
and Georgia ($-14.4$\,pp to $-7.2$\,pp).
This is the expected effect of preventing urban peeling: by forcing geographic
balance, AreaSection avoids creating maximally-efficient Republican geometries
in suburban rings.
The partisan direction is not systematic---AreaSection helps Democrats in NC/GA
but is neutral in WI and AL.

\subsubsection*{Pattern 5: NestSection and ApportionRegions Are Structurally Neutral}

Configs.~6 and 8 (ApportionRegions and NestSection) produce results very close
to config.~2 (GeoSection) on all metrics.
They are not partisan tools; they are audit and consistency tools.
Their value is not in changing partisan outcomes but in providing
structure-from-first-principles: the bisection tree is derived from mathematical
properties of $k$ (prime factorisation) or from multi-chamber consistency
requirements, not from any geographic or demographic optimisation.

\subsection{Metric Summary}

Table~\ref{tab:metric-summary} provides a cross-configuration summary aggregated
across all four states.

\begin{table}[h]
\centering
\caption{Cross-configuration metric summary, averaged across WI, NC, GA, AL.
  EC change relative to config.~2 (Geo+GeoSection = baseline).
  Gap change is arithmetic change in proportionality gap (pp).
  MM change is change in MM district count.
  County splits change relative to config.~2.}
\label{tab:metric-summary}
\small
\begin{tabular}{lrrrr}
\toprule
\textbf{Configuration} & \textbf{EC change (\%)} & \textbf{Gap change (pp)} & \textbf{MM change} & \textbf{County split change} \\
\midrule
1. Baseline & $+8.3$\est & $-5.7$\est & $-0.3$\est & $+2.5$\est \\
2. Geo+GeoSection (ref) & 0 & 0 & 0 & 0 \\
3. Geo+AreaSection & $-0.8$\est & $+4.4$\est & 0\est & $-1.3$\est \\
4. County+GeoSection & $+3.9$\est & $0$\est & 0 & $-4.8$\est \\
5. Geo+VRASection & $+3.1$\est & $+0.7$\est & $+1.0$\est & 0 \\
6. Geo+ApportionRegions & $+0.2$\pend & $0$\pend & $+0.3$\pend & 0\pend \\
7. County+AreaSection & $+6.1$\est & $+1.1$\est & 0\est & $-5.5$\est \\
8. Geo+NestSection & $+0.5$\pend & $0$\pend & $+0.3$\pend & $-1.0$\pend \\
\bottomrule
\multicolumn{5}{l}{\est\ = estimated from B-series models; \pend\ = pending B.11/B.12 implementation.}
\end{tabular}
\end{table}
